% GENERATED FILE. Edit canonical lessons, then run npm run generate:books.
% canonical-source-hash: fnv1a64:7b56153d5cb1135a
% canonical-lessons: LA-C303-furtum, LA-C303-latrocinium, LA-C303-grassator, LA-C303-praedo, LA-C303-praeda

\chapter{Fūrtum}
\label{ch:la-furtum}
\hlchaptermodality{303}

\begin{chapteropening}
\textbf{By the end of this chapter:} \emph{I can talk about theft and robbery: a theft, highway robbery, a mugger, a robber and the loot.}

Complete the last lesson of chapter 303: I can talk about theft and robbery: a theft, highway robbery, a mugger, a robber and the loot.
\end{chapteropening}

\section[fūrtum, fūrtī]{fūrtum, fūrtī — a theft}

\label{lesson:LA-C303-furtum}



\begin{quote}
Before the new one: say the Latin for to betray, then the Latin for to open up.
\end{quote}

\subsection*{You'll want to know: fūrtum, fūrtī}
\textbf{fūrtum, fūrtī} — \textquotedblleft{}a theft\textquotedblright{}.

\begin{quote}
Fūrtum in balneō factum est.
\end{quote}

\begin{quote}
\emph{A theft was committed at the baths.}
\end{quote}

Neuter, built on \emph{fūr}, 'thief', which is taught. The ablative \emph{fūrtō} means 'by stealth'.

\begin{cousinweb}
\textbf{fūrtum} is built on \emph{fūr}, \textquotedblleft{}thief\textquotedblright{} (EtymDB), which is inherited from Proto-Italic \emph{*fōr}.
\end{cousinweb}

\begin{grammarlens}[title={What you've built}]
One more word for telling, explaining and describing at B1.
\end{grammarlens}

\subsection*{Your turn}
\begin{itemize}
\raggedright
  \item \emph{Say it:} \emph{fūrtum, fūrtī}
  \item \emph{Say it:} \emph{fūrtum, fūrtī}, once more
  \item \emph{Say it:} a sentence of your own with \emph{fūrtum, fūrtī}
  \item \emph{Recall:} say the Latin for to betray, then the Latin for to open up, then say \emph{fūrtum, fūrtī} again
\end{itemize}

\begin{tcolorbox}[breakable,colback=teal!4,colframe=teal!35!black,title={Before you move on}]
What does fūrtum, fūrtī mean? (\textquotedblleft{}A theft\textquotedblright{}.) Say it once more.
\end{tcolorbox}
\section[latrōcinium, latrōciniī]{latrōcinium, latrōciniī — highway robbery, banditry}

\label{lesson:LA-C303-latrocinium}



\begin{quote}
Before the new one: say the Latin for to open up, then the Latin for a theft.
\end{quote}

\subsection*{You'll want to know: latrōcinium, latrōciniī}
\textbf{latrōcinium, latrōciniī} — \textquotedblleft{}highway robbery, banditry\textquotedblright{}.

\begin{quote}
In illā silvā latrōcinia crēbra erant.
\end{quote}

\begin{quote}
\emph{In that forest robberies were frequent.}
\end{quote}

Neuter, from \emph{latrō} (taught), 'bandit'. Cicero calls a lawless war a \emph{latrōcinium}.

\begin{cousinweb}
\textbf{latrōcinium} is built on \emph{latrō}, \textquotedblleft{}bandit\textquotedblright{}, first \textquotedblleft{}a hired soldier\textquotedblright{} (EtymDB, Lewis \& Short). EtymDB traces \emph{latrō} to Greek \emph{latron}, \textquotedblleft{}pay\textquotedblright{}.
\end{cousinweb}

\begin{grammarlens}[title={What you've built}]
One more word for telling, explaining and describing at B1.
\end{grammarlens}

\subsection*{Your turn}
\begin{itemize}
\raggedright
  \item \emph{Say it:} \emph{latrōcinium, latrōciniī}
  \item \emph{Say it:} \emph{latrōcinium, latrōciniī}, once more
  \item \emph{Say it:} a sentence of your own with \emph{latrōcinium, latrōciniī}
  \item \emph{Recall:} say the Latin for to open up, then the Latin for a theft, then say \emph{latrōcinium, latrōciniī} again
\end{itemize}

\begin{tcolorbox}[breakable,colback=teal!4,colframe=teal!35!black,title={Before you move on}]
What does latrōcinium, latrōciniī mean? (\textquotedblleft{}Highway robbery, banditry\textquotedblright{}.) Say it once more.
\end{tcolorbox}
\section[grassātor, grassātōris]{grassātor, grassātōris — a mugger, a street thug}

\label{lesson:LA-C303-grassator}



\begin{quote}
Before the new one: say the Latin for a theft, then the Latin for highway robbery.
\end{quote}

\subsection*{You'll want to know: grassātor, grassātōris}
\textbf{grassātor, grassātōris} — \textquotedblleft{}a mugger, a street thug\textquotedblright{}.

\begin{quote}
Noctū grassātōrēs per viās errābant.
\end{quote}

\begin{quote}
\emph{At night muggers roamed the streets.}
\end{quote}

Masculine, 3rd declension, from \emph{grassor}, 'to prowl'. Satirists warn of them in dark city streets.

\begin{cousinweb}
\textbf{grassātor} is built on \emph{grassor}, \textquotedblleft{}to go about, to prowl\textquotedblright{} (EtymDB), with the ending \emph{-tor}. Lewis \& Short note an older sense, \textquotedblleft{}an idler\textquotedblright{}.
\end{cousinweb}

\begin{grammarlens}[title={What you've built}]
One more word for telling, explaining and describing at B1.
\end{grammarlens}

\subsection*{Your turn}
\begin{itemize}
\raggedright
  \item \emph{Say it:} \emph{grassātor, grassātōris}
  \item \emph{Say it:} \emph{grassātor, grassātōris}, once more
  \item \emph{Say it:} a sentence of your own with \emph{grassātor, grassātōris}
  \item \emph{Recall:} say the Latin for a theft, then the Latin for highway robbery, then say \emph{grassātor, grassātōris} again
\end{itemize}

\begin{tcolorbox}[breakable,colback=teal!4,colframe=teal!35!black,title={Before you move on}]
What does grassātor, grassātōris mean? (\textquotedblleft{}A mugger, a street thug\textquotedblright{}.) Say it once more.
\end{tcolorbox}
\section[praedō, praedōnis]{praedō, praedōnis — a robber, a pirate}

\label{lesson:LA-C303-praedo}



\begin{quote}
Before the new one: say the Latin for highway robbery, then the Latin for a mugger.
\end{quote}

\subsection*{You'll want to know: praedō, praedōnis}
\textbf{praedō, praedōnis} — \textquotedblleft{}a robber, a pirate\textquotedblright{}.

\begin{quote}
Praedōnēs mercātōrem cēpērunt et aurum eius retinuērunt.
\end{quote}

\begin{quote}
\emph{Robbers captured a merchant and kept his gold.}
\end{quote}

Masculine, 3rd declension. At sea a \emph{praedō} is a pirate.

\begin{cousinweb}
\textbf{praedō} is built on \emph{praeda}, \textquotedblleft{}loot, booty\textquotedblright{} (EtymDB): one who takes booty.
\end{cousinweb}

\begin{grammarlens}[title={What you've built}]
One more word for telling, explaining and describing at B1.
\end{grammarlens}

\subsection*{Your turn}
\begin{itemize}
\raggedright
  \item \emph{Say it:} \emph{praedō, praedōnis}
  \item \emph{Say it:} \emph{praedō, praedōnis}, once more
  \item \emph{Say it:} a sentence of your own with \emph{praedō, praedōnis}
  \item \emph{Recall:} say the Latin for highway robbery, then the Latin for a mugger, then say \emph{praedō, praedōnis} again
\end{itemize}

\begin{tcolorbox}[breakable,colback=teal!4,colframe=teal!35!black,title={Before you move on}]
What does praedō, praedōnis mean? (\textquotedblleft{}A robber, a pirate\textquotedblright{}.) Say it once more.
\end{tcolorbox}
\section[praeda, praedae]{praeda, praedae — loot, booty; prey}

\label{lesson:LA-C303-praeda}



\begin{quote}
Before the new one: say the Latin for a mugger, then the Latin for a robber.
\end{quote}

\subsection*{You'll want to know: praeda, praedae}
\textbf{praeda, praedae} — \textquotedblleft{}loot, booty; prey\textquotedblright{}.

\begin{quote}
Latrōnēs praedam in spēluncā posuērunt.
\end{quote}

\begin{quote}
\emph{The bandits put the loot in a cave.}
\end{quote}

Feminine. Also the prey of an animal. English \emph{prey} comes from it through French, and \emph{predator} from its verb.

\begin{cousinweb}
\textbf{praeda} is explained by Lewis \& Short as a short form of an older \emph{*praehenda}, \textquotedblleft{}a thing seized\textquotedblright{}, from \emph{prehendō}, \textquotedblleft{}to seize\textquotedblright{}. English \emph{prey} comes from it through French.
\end{cousinweb}

\begin{grammarlens}[title={What you've built}]
One more word for telling, explaining and describing at B1.
\end{grammarlens}

\subsection*{Your turn}
\begin{itemize}
\raggedright
  \item \emph{Say it:} \emph{praeda, praedae}
  \item \emph{Say it:} \emph{praeda, praedae}, once more
  \item \emph{Say it:} a sentence of your own with \emph{praeda, praedae}
  \item \emph{Recall:} say the Latin for a mugger, then the Latin for a robber, then say \emph{praeda, praedae} again
\end{itemize}

\begin{tcolorbox}[breakable,colback=teal!4,colframe=teal!35!black,title={Before you move on}]
What does praeda, praedae mean? (\textquotedblleft{}Loot, booty; prey\textquotedblright{}.) Say it once more.
\end{tcolorbox}
